%% kapa: exam=1 round=19 year=2008 period=1 form=A subject=civil_law subject_ko=민법 seq=29 number=29
%% source: Q-Net 감정평가사 기출문제 (cst00309) artlSeq=1001962 file=extracted/1008/1001962/2008년 제19회 감정평가사 1차시험 기출문제/1교시.hwp
%% conversion: kordoc 4.24.51 -> tools/gen1.py (auto); review=proofread-sonnet; flags=none
\begin{question}{29}
甲명의로 등기되어 있는 X토지에 대해 乙이 점유취득시효의 요건을 모두 갖추었다. 다음 설명 중 틀린 것은? (다툼이 있는 경우에는 판례에 의함)
\begin{choices}
  \item 乙은 甲에게 소유권이전등기를 청구할 수 있고, 그 청구권의 법적 성질은 채권적 청구권이다.
  \item 甲이 취득시효가 완성된 사실을 알고 X토지를 제3자 丙에게 양도하고 소유권이전등기를 넘겨줌으로써 乙에 대한 소유권이전등기의무를 이행불능에 빠지게 한 경우, 甲은 乙에게 불법행위책임을 진다.
  \item 乙이 소유권이전등기를 하기 전에 甲이 X토지를 제3자 丙에게 양도하고 이전등기를 해 준 경우, 乙은 甲에게 채무불이행책임을 물을 수 없다.
  \item 乙이 소유권이전등기를 하기 전에 甲이 X토지를 제3자 丙에게 양도하고 이전등기를 해 준 경우, 丙이 乙의 취득시효사실을 알면서 소유권을 취득하였더라도 특별한 사정이 없는 한 甲의 乙에 대한 소유권이전등기의무를 丙이 승계하는 것은 아니다.
  \item 乙이 소유권이전등기를 하기 전에 甲이 X토지를 제3자 丙에게 양도하고 이전등기를 해 준 경우, 丙이 乙의 취득시효사실을 알고서 X토지를 양도받은 경우라면 乙은 丙에게 시효취득을 주장할 수 있다.
\end{choices}
\end{question}
\answer{5}
